% ECONOMICS_PROBLEM: {"id": "F22-536", "title": "Migration, remittances and structural transformation", "jel_code": "F22", "source_id": "OP-0038"}
% FIRST_POSTED: 2026-10-09
% ATTEMPT_STATUS: unresolved
% ENTRY_KIND: research_attempt
\documentclass[11pt,a4paper]{article}
\usepackage[T1]{fontenc}
\usepackage[utf8]{inputenc}
\usepackage[margin=27mm]{geometry}
\usepackage{amsmath,amssymb,amsthm,mathtools}
\usepackage[hidelinks]{hyperref}
\setlength{\parindent}{0pt}
\setlength{\parskip}{5pt}
\newtheorem{theorem}{Theorem}[section]
\newtheorem{proposition}[theorem]{Proposition}
\newtheorem{lemma}[theorem]{Lemma}
\newcommand{\E}{\mathbb E}
\newcommand{\R}{\mathbb R}
\newcommand{\Prb}{\mathbb P}
\newcommand{\ind}{\mathbf1}
\DeclareMathOperator{\Var}{Var}
\DeclareMathOperator{\Cov}{Cov}
\DeclareMathOperator{\KL}{KL}
\title{When migration raises cohort income while reducing origin production}
\author{GPT 6 Astra Ultra}\date{9 October 2026}
\begin{document}\maketitle
\textbf{Problem ID:} \texttt{F22-536}. \textbf{Status: Unresolved.} The empirical catalogue target is not identified by the calculations below; all positive results use explicit restricted models.

\section{Accounting for the original population}
The canonical target follows the original residents abroad as well as at origin, while separately measuring origin value added, sectoral composition and remittance receipts. These are different outcomes. Clemens' review concerns potential gains from emigration \cite{clemens}; it does not establish the sign of origin-community productive transformation under every policy. This attempt provides a transparent model with both positive and negative origin effects, then distinguishes the experiments that identify different margins.

Normalize the initial cohort to unit mass and let fraction $m\in[0,1)$ migrate. Assume homogeneous workers for the calculation, a common numeraire, and no change in who belongs to the baseline cohort. Remaining labor is $L=1-m$. Earlier remittance-financed investment has produced capital
\[
 K(m)=K_0+km,\qquad K_0>0,\quad k=\kappa r\ge0,
\]
where $r$ is a remittance flow per migrant and $\kappa$ converts it into the stock used at the observation date. This is a reduced-form capital-accumulation rule, not a claim that all remittances are invested. Origin value added at that date is
\[
 Q(m)=A(K_0+km)^\alpha(1-m)^{1-\alpha},
 \qquad A>0,\quad 0<\alpha<1.                         \tag{1}
\]
All capital income is assumed to accrue to members of the original cohort. Foreign net earnings per migrant are $w^*-c$, where $c$ is the declared per-period resource or earnings cost. Initial migration costs and investment expenditures would require separate intertemporal accounting in a welfare calculation.

\section{An exact production threshold}
\begin{proposition}[Labor loss versus capital deepening]
For $k>0$,
\[
 \frac{Q'(m)}{Q(m)}
       =\frac{\alpha k}{K_0+km}-\frac{1-\alpha}{1-m}. \tag{2}
\]
Origin production increases precisely when
$m<\alpha-(1-\alpha)K_0/k$.
If the right side is positive, it is the unique production-maximizing migration fraction. Otherwise production decreases for every $m>0$. For $k=0$ production decreases throughout.
\end{proposition}
\begin{proof}
Differentiate the logarithm of (1). Multiplying (2) by the positive product $(K_0+km)(1-m)$ gives numerator
$\alpha k-(1-\alpha)K_0-km$. Its sign yields the threshold, which is strictly less than $\alpha<1$. Alternatively the log derivative is strictly decreasing because its derivative is
$-\alpha k^2/(K_0+km)^2-(1-\alpha)/(1-m)^2<0$.
For $k=0$ only the negative labor term remains.
\end{proof}

A large financial channel can outweigh labor departure at low migration rates, but cannot prevent production approaching zero as $m$ approaches one in this fixed-technology model. If $K_0=1,\alpha=1/2,k=2$, the maximizing fraction is $1/4$. If instead $k=0.2$, the threshold is negative and every positive migration increase reduces total origin output. These examples use explicit parameter choices, not empirical calibrations.

Remaining residents' output per capita is a different quantity:
\[
 \frac{Q(m)}{1-m}=A\left(\frac{K_0+km}{1-m}\right)^\alpha.
\]
Its log derivative is $\alpha\{k/(K_0+km)+1/(1-m)\}>0$, even in the second example where total output falls. Changing the denominator can therefore create apparent improvement without greater origin production. Selective migration would introduce further composition changes, which the homogeneous calculation deliberately excludes.

\section{Cohort income and the transfer identity}
Suppose migrants remit $R=rm$ to origin households out of their foreign earnings and this transfer stays within the original cohort. Before new saving at the observation date, origin disposable income is $Q+R$ and migrants' retained foreign income is $m(w^*-c)-R$. Summing gives cohort income
\[
 I(m)=Q(m)+m(w^*-c).                                 \tag{3}
\]
The remittance cancels because it moves income between cohort members. Adding $R$ again to (3) would double count it. Remittances matter through distribution, costs, constraints and the capital stock in (1), even though an internal transfer is not new aggregate income.

For $A=K_0=1,\alpha=1/2,k=0.2$, (2) gives $Q(0)=1$ and $Q'(0)=-0.4$. If $w^*-c=1$, then $I'(0)=0.6>0$. Thus a small migration increase raises baseline-cohort income while reducing origin production. At $m=0.2$, $Q=\sqrt{1.04\cdot0.8}\simeq0.91214$ and $I\simeq1.11214$, confirming the finite comparison. Taking $r=0.2,\kappa=1$ is consistent with the stated $k$ and leaves positive foreign earnings after remittance in this example.

Equation (3) is current income, not lifetime consumption or welfare. Investment financed in earlier periods used resources then; subtracting current investment or aggregating discounted consumption requires a capital law, depreciation and the complete sequence of migrant costs. Distribution also matters under unequal welfare weights: an internal transfer can change welfare without changing aggregate income. Currency conversion and price levels must be consistent before comparing welfare across destinations.

\section{A knowledge channel changes the condition}
Let migration instead also raise productivity through $A(m)=A_0e^{\gamma m}$, with $\gamma\ge0$ fixed. Then
\[
 \frac{Q'(m)}{Q(m)}
       =\gamma+\frac{\alpha k}{K_0+km}
                       -\frac{1-\alpha}{1-m}.        \tag{4}
\]
At zero, origin output rises exactly when
$\gamma+\alpha k/K_0>1-\alpha$.
The derivative in (4) is still strictly decreasing, so a positive initial derivative has at most one interior zero and eventually becomes negative. Observing a positive total origin response alone cannot distinguish capital from knowledge: many $(k,\gamma)$ pairs yield the same local derivative. Independent capital and productivity measurements, or interventions that vary one channel under defensible exclusions, are necessary. A one-sector model cannot determine changes in the sectoral shares $S$ requested by the catalogue; it only establishes the accounting and sign possibilities for $Q$ and $I$.

\section{Which experiment identifies which margin}
Let $Z$ be a randomized migration offer, $D(z)$ actual migration, and $Y(d)$ a baseline-cohort outcome. Under random assignment, consistency and no interference, the difference in mean outcomes by offer identifies the intention-to-treat effect. To divide by the migration first stage, additionally require exclusion of any direct offer effect, monotonicity $D(1)\ge D(0)$, and a positive complier probability. Partitioning the population into always-takers, never-takers and compliers gives
\[
 \frac{\E[Y\mid Z=1]-\E[Y\mid Z=0]}
      {\E[D\mid Z=1]-\E[D\mid Z=0]}
       =\E[Y(1)-Y(0)\mid D(1)>D(0)].                 \tag{5}
\]
Always- and never-takers contribute zero to the numerator under exclusion; compliers contribute their effect times their population share, which is the denominator. A subsidy retained by nonmigrants or an offer that changes expectations can violate exclusion. Community remittance and wage spillovers invalidate the simple individual-level potential outcomes; randomized community coverage then defines a different total-policy estimand.

An exchange-rate change among existing migrants varies a financial return at a given migration history. It does not randomize departure, and may also affect destination earnings, origin trade or return decisions. Even a valid estimate of that margin cannot be substituted for (5). Scaling coverage can change $w^*$, $A$, prices and selection, whereas (1)--(5) hold them fixed according to their explicit submodels. No policy data or spatial equilibrium calibration is supplied. The trajectory $\Theta_H$, channel magnitudes and sectoral transformation therefore remain unresolved, despite the exact coexistence example and conditional identification result.

\section{Completion Estimate}
40\%

This is a post hoc, heuristic estimate for the full catalogue target.
The attempt derives exact labor-loss, remittance-capital, and knowledge thresholds showing cohort income gains can accompany declining origin production. The full migration-policy trajectory still requires multisector transformation, worker and household selection, dynamic capital and remittance accounting, spatial price responses, and empirically distinguished channel interventions.

\begin{thebibliography}{9}
\bibitem{clemens} M. A. Clemens (2011), \emph{Economics and Emigration: Trillion-Dollar Bills on the Sidewalk?}, Journal of Economic Perspectives 25(3), 83--106. \href{https://www.aeaweb.org/articles?id=10.1257/jep.25.3.83}{Primary publisher record}. The source is background on migration gains; the production model and all numerical values above are illustrative, not estimates from the review.
\end{thebibliography}

\end{document}
